\documentclass[11pt, oneside]{article} 
\usepackage{geometry}
\geometry{letterpaper} 
\usepackage{graphicx}
	
\usepackage{amssymb}
\usepackage{amsmath}
\usepackage{parskip}
\usepackage{color}
\usepackage{hyperref}

\graphicspath{{/Users/telliott/Dropbox/Github-math/figures/}}
% \begin{center} \includegraphics [scale=0.4] {gauss3.png} \end{center}

\title{Polar area problem}
\date{}

\begin{document}
\maketitle
\Large

%[my-super-duper-separator]

Here is a problem from the web:
\begin{center} \includegraphics [scale=0.4] {circ_seg_prob.png} \end{center}
The difficult part is the area in the upper-right hand corner.  We need to find the area of the red segment that is cut by the angled line (or its complement in white).

Let's focus on the area of that part of a circle swept out by a ray from one edge, i.e. between a chord and the diagonal.

\begin{center} \includegraphics [scale=0.4] {polar_area.png} \end{center}
The figure was constructed by drawing a circle inside a square, with the horizontal diagonal of the circle bisecting the square, and then only the top half of the square and circle are shown.

Connect the left end of the diagonal with the edge of the square on the opposite side.  The angle $\theta$ formed has the tangent $1/2$.  We wish to calculate the area enclosed by the diagonal, the ray forming angle $\theta$ with it, and the edge of the circle.

Since $r$ varies with $\theta$, we need an expression for $r(\theta)$.

\subsection*{geometry}
We don't actually need calculus for this problem.  Drawing lines from the center of the circle, the area can be divided into two right triangles with sides $R \sin \theta$ and $a = R \cos \theta$, plus a sector of the circle of radius $R$ swept out by angle $\phi = 2 \theta$.

\begin{center} \includegraphics [scale=0.4] {polar_area.png} \end{center}
To make things a bit simpler, we will assume now that this is a unit circle with $R = 1$, recognizing that the result can always re-scaled.  For a circle of radius $R$, the area scales like $R^2$.

The area of the each triangle is $a \sin \theta / 2$, so the two together are $A = \sin \theta \cos \theta$.  Since this is a unit circle, the area of a sector with angle $\phi$ is:
\[ \frac{A}{\pi} = \frac{\phi}{2 \pi} = \frac{2 \theta}{2 \pi} \]
This follows from the inscribed angle theorem, which gives $\phi = 2 \theta$.

So the sector is $\theta$ and the total is
\[ A = \theta + \sin \theta \cos \theta \]

In this particular case, we can do even better.  We have that $\tan \theta = 1/2$, meaning that $2 \sin \theta = \cos \theta$, and then from the Pythagorean theorem $5 \sin^2 \theta = 1$, so $\sin \theta = 1/\sqrt{5}$ and
\[ A = \theta + \frac{2}{5} \]

\subsection*{calculus:  setup}

If the length of the ray is $r$ then the "sides" of the relevant area element are $r$ and $r \ d \theta$, and the area is
\[ A = \frac{1}{2} \  \int_0^{\theta} r \cdot r \ d \theta =  \frac{1}{2} \  \int_0^{\theta} r^2 \ d \theta \]

One way to get an expression for $r(\theta)$ is to consider the triangle with two sides that are radii of the circle, and third side $r$.  We have that $a = R \cos \theta$ and $r = 2a$ so $r = 2 R \cos \theta$.  

This fits with what we already know from Thales' theorem.  The triangle with sides $r$ and $2R$ is a right triangle.

We might also use the law of cosines
\[ R^2 = r^2 + R^2 - 2rR \cos \theta \]
\begin{center} \includegraphics [scale=0.4] {polar_area.png} \end{center}

This gives a quadratic in $r$
\[ r^2 = 2R \cos \theta \ r \]
whose solution is
\[ r = 2R \cos \theta \]

For the unit circle we have just 
\[ r = 2 \cos \theta \ \ \ \ \ \  r^2 = 4 \cos^2 \theta \]

It's a bit of a diversion, but a different approach is to consider the central angle $\phi = 2 \theta$.  The supplementary angle to $\phi$ (call it $\phi'$) is the apex angle in the same double triangle we had before.  Its cosine is the same magnitude as $\cos \phi$ but with a minus sign.

The law of cosines gives here:
\[ r^2 = R^2 + R^2 - 2 R^2 \cdot (- \cos 2 \theta) \]
For the unit circle
\[ r^2 = 2 + 2 \cos 2 \theta \]

Compared to what we had before it seems a little different
\[ r^2 = 2 + 2 \cos 2 \theta \stackrel{?}{=} 4 \cos^2 \theta \]
But this is the double angle formula in disguise.  Rearranging:
\[ 1 + \cos 2 \theta = 2 \cos^2 \theta = \cos^2 \theta + 1 - \sin^2 \theta \]
\[ \cos 2 \theta = \cos^2 \theta - \sin^2 \theta \]
which is correct.

\subsection*{calculus:  integral}

The integral is 
\[ A = \frac{1}{2} \ \int_0^{\theta} r^2 \ d \theta \]
Using the first version this is 
\[ A = 2 \int \cos^2 \theta \ d \theta = \theta + \sin \theta \cos \theta \bigg |_0^{\theta} \]
\[ = \theta + \sin \theta \cos \theta \]

and using the second version it is
\[ A = \frac{1}{2} \ \int_0^{\theta}  2 + 2 \cos 2 \theta \ d \theta \]
\[ = \int_0^{\theta}  1 + \cos 2 \theta \ d \theta \]
\[ = \theta + \frac{1}{2} \ \sin 2 \theta \ \bigg |_0^{\theta} = \theta + \sin \theta \cos \theta \]

\subsection*{displaced circle}

The equation for a unit circle of radius $R$ centered at the origin is simply 
\[ R^2 = x^2 + y^2 \]
Every $(x,y)$ on the circle satisfies this equation.

For a unit circle displaced from the origin to $(1,0)$, the equation is
\[ 1^2 = (x-1)^2 + y^2 \]
This corresponds with the origin of coordinates at the bottom-left corner of the figure.
\begin{center} \includegraphics [scale=0.4] {polar_area.png} \end{center}

Expanding
\[ 1 = x^2 - 2x + 1 + y^2 \]

For the ray $r$, $x = r \cos \theta$ and $y = r \sin \theta$ so $x^2 + y^2 = r^2$ and then
\[ 1 = r^2 - 2x + 1 \]
Substituting for $x$
\[ r^2 = 2 r \cos \theta \]
\[ r = 2 \cos \theta \]
which matches what we had from the law of cosines.

\subsection*{yet another view}
Another way to look at this problem is to \emph{turn the circle}.  In the figure below, $r$ has been turned so that it is horizontal.  Part of the area is beneath the horizontal diagonal of the circle, but there is a matching sector above the diagonal, on the left.
\begin{center} \includegraphics [scale=0.4] {polar_area2.png} \end{center}
The area we need is that between $r$ and the diagonal.

The equation of the unit circle as a function $y = f(x) = \sqrt{1- x^2}$.

Consider a chord of the unit circle, drawn at a height $h$ above the center (measured on the chord's perpendicular bisector).  We must find the two points where the chord meets the circle.  Let's call them $\pm \ x_0$.  These are the two endpoints which are on the circle and where the corresponding $y$-value is $h$.

\[ y = f(x_0) = \sqrt{1 - x_0^2} = h \]
Rearranging
\[ x_0 = \pm \ \sqrt{1 - h^2} \]

The total area is
\[ A = \int_{-x_0}^{x_0} \sqrt{1 - x^2} \ dx \]

This integral is famous.  If we do a trig substitution where $x = \sin \phi$ so then $\sqrt{1 - x^2} = \cos \phi$, and also $dx = \cos \phi \ d \phi$, the integral is 
\[ \int \cos^2 \phi \ d \phi = \frac{1}{2} (\phi + \sin \phi \cos \phi) \]
which we solved without comment above.

It is easily checked:
\[ \frac{d}{d \phi} \ \frac{1}{2} (\phi + \sin \phi \cos \phi) = \frac{1}{2} (1 - \sin^2 \phi + \cos^2 \phi) \]
\[ = \frac{1}{2} (2 \cos^2 \phi) = \cos^2 \phi \] 

We could integrate the area of the circle between $-x_0$ and $+x_0$, find the part which lies below the line $y = h$, and then add the pieces from the ends.  There is a better way.

\subsection*{circular segment}
Having rotated the area, we can see that this problem can also be solved using only geometry in the new view:
\begin{center} \includegraphics [scale=0.3] {circ_seg1b.png} \end{center}
Since $r$ is horizontal, we have some equal angles $\theta$ by alternate interior angles.  The area of each of those wedges is $\theta/2$.

We calculate the total area of the two central triangles joined together as .$\sin \theta \cos \theta$.  This is looking familiar.  

So the area of the region below the circular ``cap" is just $\theta + \sin \theta \cos \theta$, as we've been saying.  The cap itself is $\pi/2$ minus this.

Another formulation for the region below the cap is $\theta + \frac{1}{2} \sin 2 \theta$, so the cap itself is
\[ A_{\text{cap}} = \frac{1}{2} ( \pi - 2 \theta - \sin 2 \theta) \]
This is not really any different than what we did before but it does solve the cap problem.

\subsection*{vertical view}
The last view (which is so simple that I'll skip the diagram for a moment), is to turn the circle another quarter-turn.  The area that we want is simply
\[ A = \int_0^h  \sqrt{1 - x^2} \ dx \]
We need a further factor of $2$ for the final answer because this is only the area \emph{above} the $x$-axis.

In terms of the trig substitution (and multiplying by that factor of $2$)  this is
\[ A= \phi + \sin \phi \cos \phi \]
Switching back to the variable $x$ we have
\[ A = \sin^{-1} x + x \ \sqrt{1 - x^2} \]
The first term on the right-hand side is the inverse sine, which can be read as ``the angle whose sine is $x$".

The bounds are simply $0$ and $h$.  But since $x=0$ at the lower bound, the expression evaluates to $0$ there.  At the upper bound we have
\[ A =  \sin^{-1} h + h \ \sqrt{1 - h^2} \]

What is $h$?  If you look back to the first diagram you should be able to see that $h = \sin \theta$, so finally
\[ A =  \theta + \sin \theta \cos \theta \]

This is equal to the answers from other approaches, as we should require.

The other nice thing about this is we see again where the terms come from --- as in the geometric approach above.  The area under the curve can be constructed in two parts by drawing the radius to the point $(x,y)$.

\begin{center} \includegraphics [scale=0.4] {circle_integral.png} \end{center}
There is a sector of the circle and a triangle.  The area is
\[ A = \frac{1}{2} \ (\sin^{-1} x + x \sqrt{1-x^2} ) \]

Accounting for the duplicated area below the $x$-axis, these correspond to the terms of
\[ A = \sin^{-1} x + x \sqrt{1-x^2} \]

It is easy to see how this relates to $\theta$:
\[ A = \theta + \sin \theta \cos \theta \]

\subsection*{cartesian coordinates}

Let's try to calculate the same area with the original orientation and without displacing the circle, i.e. with the origin of coordinates at the center of the circle.
\begin{center} \includegraphics [scale=0.4] {polar_area3.png} \end{center}

We will call the point where the ray intercepts the circle $(x_0,y_0)$.  We can then break the figure up into a triangle with base $1 + x_0$, height $y_0$, and area
\[ A_{\triangle} = \frac{1}{2} \ (1 + x_0) y_0 = \frac{y_0}{2} + \frac{x_0 y_0}{2} \]

We note that 
\[ y_0 = \sin \phi = \sin 2 \theta = 2 \sin \theta \cos \theta \]
Hence
\[ A_{\triangle} = \sin \theta \cos \theta + \frac{x_0 y_0}{2} \]
So that is one of the terms that we need to match the other answers.  

Next, the tip of the circle is (using the integral we solved earlier):
\[ A_{\odot} = \frac{1}{2} (\sin^{-1} x + x \ \sqrt{1 - x^2} \bigg |_{x_0}^1 ) \]

Neglect the leading factor of $1/2$ for the moment.  At the upper bound the first term is $\pi/2$, because the sine of $\pi/2$ is equal to $1$.  

Obviously, we have to get rid of that, but how?

Considering the lower bound the first term is $\sin^{-1} x_0$.  Going back to the diagram, $x_0$ is the \emph{cosine} of the base angle $\phi$, so it is the sine of $\pi/2 - \phi$.  For the first term only, upper bound minus lower bound, and still ignoring that factor of $2$:
\[ \pi/2 - \sin^{-1} ( \sin (\pi/2 - \phi)) = \phi \]
The whole thing is then
\[ A_{\odot} = \frac{1}{2} (\phi - x_0 \sqrt{1 - x_0^2} \ ) \]
But $\phi = 2 \theta$ and the second term is also $x_0 y_0$ so this becomes
\[ A_{\odot} = \theta - \frac{x_0 y_0}{2} \]
Hence
\[ A_{\text{tot}} = A_{\triangle} + A_{\odot} = \theta + \sin \theta \cos \theta \]

$\square$

\end{document}